\newcommand{\ratio}[1]{\frac{\avgPost[{#1}][\cMass{}]}{\cMass{#1}}}
\newcommandx{\ratioFull}[3][1=\w, 2=\w, 3=\cMass{}, usedefault]{
    \frac{\avgPost[{#1}][#3]}{#3(#2)}
    \prod_i \frac{\avgPost[{\w[i]}][#3]}{#3(\w[i])}
}
\newcommandx{\pusFrontal}[2][1={\avgPost[][\cMass{}]}, 2=\cMass{}]{
    \left(\frac{#1(\w)}{#2(\w)} \prod_i \frac{#1(\w[i])}{#2(\w[i])} \right)
}
\newcommand{\fixedPoint}{\normConst \solPrior}
\newcommand{\fpMap}[1]{T \ifthenelse{\isempty{#1}}{}{\left[\,#1\,\right]}}
\newcommand{\normConst}{\mu}
\newcommand{\solPrior}{r}
\newcommand{\solDist}{R}

\newcommandx{\leg}[2][1={\w[i]}, 2=\solPrior, usedefault]{
    \frac{\avgPost[{#1}][{#2}]}{\ifthenelse{\isempty{#2}}{\unnorm(#1)}{#2(#1)}}
}
\newcommandx{\frontalPus}[2][1=\w, 2=\solPrior, usedefault]{\left( \leg[{#1}][{#2}] \prod_i \leg[{\w[i]}][{#2}] \right)}
\newcommand{\unnorm}{\rho}

In the dice example, following the principle of objectivity led us to define a query set that included vectors
representing $n$ trials of our experiment.
This query set contained two variables, \X and \Y: \X represented the sample space,
while \Y corresponded to trials of a two-sided dice indicating whether the original dice result was odd or even.

In many cases, we consider a query set structured as follows:
\[
    \Qset = \{\X, \seqRv{k}\},
\] where \X represents the whole sample space, and each \X[i] represents $n$ trials of a transformed experiment.
We usually assume a bounded, but sufficiently large number of trials.
The transformed experiments are functions of the original experiment and can be parametrized by functions of the
model parameter.
In the next theorem, for query sets of this form, we will develop a fixed point equation characterizing the
central distribution.



\begin{theorem}
    \label{th:fixed_point}
    Suppose \pmodel is a generic parametric model, and $\Qset = \{\X, \seqRv{k}\}$ is a query set in \Mt.
    Let $\W[i]$, for every $i$, be a function of the parameter $\param$ that fully parametrizes \X[i],
    That is, for $\w[i] = \W[i](\paramValue)$ we have:
    \begin{equation}
        \pi(\xIt[i]) = \pi \left( \xIw[i] \right)\cdot
        \label{eq:func_of_theta}
    \end{equation}
    Also, let $\W(\paramValue) = \vSeq{\w}{k} = \w$.

    We define a map, denoted by $\fpMap{}$, on the space of unnormalized probability measures of \w, as follows:
    \begin{equation}
        \fpMap{\unnorm}(\w) = \frontalPus[][\unnorm]^{\frac{1}{k+1}} \unnorm(\w),
        \label{eq:pus_frontal}
    \end{equation}
    where $\unnorm$ is an unnormalized probability measure, and:
    \begin{align*}
        \avgPost[{\w[i]}][\rho] &= \exp \left\{\sum_{\x[i]} \pi(\xIw[i]) \ln \rho(\wIx[i]) \right\},\\
        \avgPost[\w][\rho] &= \int_{\omegaArea} \avgPost[\paramValue][\rho] d\paramValue,\\
        \avgPost[\paramValue][\rho] &= \exp \left\{\sum_{\x} \pi(\xIt) \ln \rho(\tIx) \right\} \cdot\\
    \end{align*}

    Let $\solPrior$ be a strictly positive prior probability distribution.
    The unnormalized measure $\fixedPoint$ is a fixed point for $\fpMap{}$, and
    \begin{equation}
        \solPrior(\tIw) = \frac{\avgPost[\paramValue][\solPrior]}{\avgPost[\w][\solPrior]},
        \label{eq:pus_th_conditional}
    \end{equation}
    if and only if
    \begin{equation}
        \DxCond[\solDist] = (k + 1) \ln \normConst,
        \label{eq:pus_th_Dx_is_const}
    \end{equation}
    where $\solDist$ is the probability measure induced by $\solPrior$.
\end{theorem}

\begin{proposition}
    We have
    \begin{equation*}
        \Cx[\solPrior] = (k + 1) \ln \normConst \cdot
    \end{equation*}
\end{proposition}

\begin{proof}
    \newcommand{\someDist}{P}
    \newcommand{\somePrior}{p}
    To begin, consider that for any arbitrary probability distribution $\someDist$, induced by a strictly positive
    prior $\somePrior$, using Equation~\ref{eq:func_of_theta}, we can derive the following:
    \begin{align*}
        \DxCond[\someDist] &= \KlCond[\someDist][\rval{x}][ ] + \DxCondExpandFull[\someDist][][ ] \\
        &= \sum_{\rval{x}} \pi(\xIt) \ln \frac{\somePrior(\tIx)}{\somePrior(\paramValue)} +
        \sumix \pi(\xIw[i]) \ln \frac{\somePrior(\wIx[i])}{\somePrior(\w[i])}\\
        &= \ln \leg[\paramValue][\somePrior] + \sum_i \ln \leg[][\somePrior] \\
        &= \ln \frontalPus[\paramValue][\somePrior] \cdot
    \end{align*}
    Since $\W(\paramValue) = \w$, we have $\somePrior(\paramValue) = \somePrior(\tIw) \somePrior(\w)$, and
    \begin{equation}
        \label{eq:pus_proof_main_eq}
        \DxCond[\someDist] + \ln \somePrior(\tIw) =
        \ln \left(  \frac{\avgPost[\paramValue][\somePrior]}{\somePrior(\w)} \prod_i \leg[][\somePrior] \right) \cdot
    \end{equation}

    To prove the `if' part, we assume Equation~\ref{eq:pus_th_Dx_is_const} holds for a probability measure $\solDist$,
    induced by a strictly positive prior $\solPrior$, for every value of \paramValue.
    We exponentiate both sides of Equation~\ref{eq:pus_proof_main_eq} to obtain
    \begin{equation}
        \normConst ^{k+1} \solPrior(\tIw) =\frac{\avgPost[\paramValue][\solPrior]}{\solPrior(\w)} \prod_{i=1}^k \leg \cdot
        \label{eq:pus_proof_a}
    \end{equation}
    Since this equation holds for any value of \paramValue, we can integrate both sides with respect to \paramValue:
    \[
        \int_{\omegaArea} \normConst ^{k+1} \solPrior(\tIw) d\paramValue =
        \int_{\omegaArea}
        \left( \frac{\avgPost[\paramValue][\solPrior]}{\solPrior(\w)} \prod_i \leg  \right) d\paramValue \cdot
    \]
    That implies
    \begin{equation}
        \normConst ^{k+1} = \frontalPus \cdot
        \label{eq:pus_proof_b}
    \end{equation}
    Now we evaluate $\fpMap{}$ at $\fixedPoint$:
    \begin{align*}
        \fpMap{\fixedPoint}(\w) &= \frontalPus[][\fixedPoint] ^{\frac{1}{k+1}} \fixedPoint(\w) \\
        &= \left( \normConst^{k+1} \right)^{\frac{1}{k+1}} \solPrior(\w) \\
        &= \fixedPoint(\w) \cdot
    \end{align*}
    Therefore, $\fixedPoint$ is a fixed point for $\fpMap{}$.
    In addition, for $\solPrior(\tIw)$ we can obtain Equation~\ref{eq:pus_th_conditional} by dividing
    Equation~\ref{eq:pus_proof_a} by Equation~\ref{eq:pus_proof_b}.
    This completes the proof of the forward direction.

    We now prove the `only if' part.
    Suppose $\fixedPoint$ is a fixed point of $\fpMap{}$, where $\solPrior$ is a strictly positive probability
    distribution and $\normConst$ is a constant.
    Since $\solPrior(\w) \neq 0$, by letting $\fpMap{\fixedPoint} = \fixedPoint$, for every value of \w we get
    \begin{equation}
        \label{eq:puss_proof_pus_is_mu}
        \frontalPus^{\frac{1}{k+1}} = \normConst \cdot
    \end{equation}

    On the other hand, under the assumptions of the theorem, Equation~\ref{eq:pus_th_conditional} holds.
    Substituting this into Equation~\ref{eq:pus_proof_main_eq} gives
    \[
        \DxCond[\solDist] = \ln \frontalPus \cdot
    \]
    Now, substituting Equation~\ref{eq:puss_proof_pus_is_mu} into the preceding equation yields
    \[
        \DxCond[\solDist] = (k + 1) \ln \normConst,
    \]
    which concludes the proof.
\end{proof}


Using Equation~\ref{eq:pus_frontal}, we can construct a sequence of unnormalized measures by
\[\rho_{n+1} = \fpMap{\rho_n} \cdot \]
This sequence can be used in an iterative algorithm to find an estimation of the central prior $\cMass{\w}$.
Once $\cMass{\w}$ is determined, Equation~\ref{eq:pus_th_conditional} can be used to obtain $\cMass{\paramValue{}}$.

The key property of the mapping $\fpMap{}$ is that it is defined on the space of unnormalized measures.
This allows us to compute $\cMass{\paramValue}$ for
specific regions of the parameter space without needing to scan the entire space, which is typically large and intractable.
Additionally, it works with \w and its marginals \w[i], rather than \paramValue.
Notably, the space of \w is significantly smaller than that of \paramValue.

In Equation~\ref{eq:pus_frontal}, if we replace $\avgPost[][\rho]$ with $\avgPost[][\cMass{}]$, it can improve the
convergence properties considerably.
In models that the posterior of the parameter converges, we can calculate $\avgPost[][\cMass{}]$ by using any well
behaved prior instead of $\cMass{}$.
We typically use a conjugate prior to simplify the derivation of the posterior potential.
By using a conjugate prior, the posterior potential takes a form similar to the entropy of the likelihood distribution.


\begin{definition}[Posterior Potential]
    Let \pmodel be a parametric model, and \X be a random variable defined on \set{s}.
    Let $\W$ be a function of the parameter \param that fully parametrizes \X.

    If the unnormalized measure $\avgPost[\w]$, defined as
    \begin{equation*}
        \avgPost[\w] = \avgPostExpand[p]{\w}{\x},
    \end{equation*}
    dose not depend on $p(\w)$, we refer to it as the \emph{posterior potential} of \w.
    \emph{(This is not formal enough)}
\end{definition}

\begin{theorem}
    Let $\fpMap{}$ be the map defined as
    \begin{equation*}
        \fpMap{\unnorm}(\w) = \frontalPus[][ ]^{\frac{1}{k+1}} \unnorm(\w),
    \end{equation*}
    where $\avgPost[\cdot]$ denotes the posterior potential.
    The sequence \[\rho_{n+1} = \fpMap{\rho_n},\] converges to the fixed point of
    $\fpMap{}$, for any arbitrary initial unnormalized probability measure $\rho_1$.
\end{theorem}

\begin{proof}
    \renewcommand{\d}[2]{d_{\W} \ifthenelse{\isempty{#1}}{}{\left( #1, #2 \right)}}
    \newcommand{\dExpand}[2]{\intAbs{\lnfrac{#1}{#2}}}
    \newcommand{\lnfrac}[2]{\ln \frac{#1(\w)}{#2(\w)}}
    \newcommand{\intAbs}[1]{\int \left| #1 \right| d\w}
    \newcommand{\p}{\rho}
    \newcommand{\q}{\eta}
    \renewcommand{\r}{\phi}
    We prove the theorem by using Banach fixed point theorem~\cite{Agarwal2018}.
    Let $\p$ and $\q$ be two unnormalized probability measures of a random variable \w
    We define the distance of $\p$ and $\q$ as follows:
    \begin{equation}
        \d{\p}{\q} = \dExpand{\p}{\q} \cdot
        \label{eq:d_metric}
    \end{equation}
    Obviously, if $\p = \q$, then $\d{\p}{\q} = 0$.
    Conversely, if $\d{\p}{\q} = 0$, that implies  $\p = \q$ because $\intAbs{f(\w)} = 0$ if and only if $f(\w) = 0$.
    Moreover, $\d{}{}$ is symmetric $\d{\p}{\q} = \d{\q}{\p}$, and satisfies the triangle inequality:
    \begin{align*}
        \d{\p}{\r} + \d{\r}{\q} &= \dExpand{\p}{\r} + \dExpand{\r}{\q} \\
        &= \int \left( \left| \lnfrac{\p}{\r} \right| + \left| \lnfrac{\r}{\q} \right| \right) d\w \\
        &\geq \intAbs{\lnfrac{\p}{\r} + \lnfrac{\r}{\q}} \\
        &= \d{\p}{\q} \cdot
    \end{align*}
    Thus, $\d{}{}$ is a metric.

    Now, we show that $\fpMap{}$ is a contraction.
    \begin{align*}
        \d{\fpMap{\p}}{\fpMap{\q}} &= \dExpand{\fpMap{\p}}{\fpMap{\q}} \\
        &= \intAbs{ \ln \left( \frac{\p(\w)}{\q(\w)} \pusFrontal[\q][\p]^{\frac{1}{k+1}} \right) }\\
        &= \frac{1}{k+1} \intAbs{ \ln \left( \prod_i \frac{p_{\p}(\w \mid \w[i])}{p_{\q}(\w \mid \w[i])} \right)} \\
        &\leq \frac{1}{k+1} \sum_{i=1}^{k} \intAbs{\ln \frac{p_{\p}(\w \mid \w[i])}{p_{\q}(\w \mid \w[i])}} \\
        &\leq^{?} \frac{k}{k+1} \d{\p}{\q}
    \end{align*}
    \emph{Unfortunately, this proof is not complete.}
\end{proof}
